\begin{annexitem}{annex:p15.pq.validities}{thm:p15-pq-validities}{Counting semantics and actual countermodels}
\begin{proof}
Fix a structure, parameters, and \(N=n^\ell>0\).
Let \(B\subseteq A^\ell\) be the scope's truth set. Pointwise implication
gives inclusion of truth sets, hence preserves a lower bound on
\(|B|/N\). Decreasing a threshold also preserves that bound.
Universal truth makes \(|B|=N\). Conversely \(|B|/N\ge r>0\)
makes \(B\) nonempty, giving the ordinary existential witness.
If both complementary scopes met thresholds \(r,s\) with \(r+s>1\),
their disjoint cardinalities would give \(N\ge(r+s)N>N\), a contradiction.
These are validity statements for all assignments, not inference from
truth of an unquantified scope at just one tuple.

In the three-element relation displayed in the body, row sizes are
\(2,2,0\), and column sizes are \(2,1,1\). The two nested evaluations
therefore have outer truth fractions \(2/3\) and \(1/3\).
For the four-element unary example each complementary truth fraction
is exactly \(1/2\), so both non-strict half-threshold assertions hold.
In general negating a threshold means \(|B|/N<r\), equivalently
the complement has fraction strictly greater than \(1-r\), not merely
at least \(1-r\). The strict inequality is a semantic description here,
not a new native binder attributed to the source.
At threshold zero the implication to ordinary existence would fail on
an empty truth set; that excluded endpoint is not silently introduced.
\end{proof}
\end{annexitem}

\begin{annexitem}{annex:p15.pq.translation}{thm:p15-pq-translation}{The full witness translation at a fixed universe size}
\begin{proof}
Translate atoms identically, commute with Boolean connectives, and
translate the scopes of ordinary quantifiers, renaming bound variables
freshly as necessary. At a probability binder put \(N=n^\ell\) and
\(k=\lceil rN\rceil\). Choose \(k\) fresh variable tuples
\(\bar z^1,\ldots,\bar z^k\) of length \(\ell\), disjoint from the
free parameters and from all scope binders after fresh renaming.
Replace the binder by
\[
 \exists\bar z^1\cdots\exists\bar z^k
 \left[
  \bigwedge_{j=1}^k\tau_n(\varphi)(\bar x,\bar z^j)
  \land
  \bigwedge_{i<j}\ \bigvee_{t=1}^{\ell}z_t^i\ne z_t^j
 \right].
\]
The second conjunction demands distinct tuples, not pairwise distinct
coordinates within one tuple. Thus repeated tuple values are counted
correctly. Since \(0<r<1\), \(1\le k\le N\).
If the probability formula is true, its finite truth set has at least
\(k\) elements, so enumerate any \(k\) distinct tuples to satisfy this
witness formula. Conversely witnesses give \(k\) distinct truth tuples,
and hence \(|B|/N\ge k/N\ge r\). The equivalence
\(|B|\ge rN\) iff \(|B|\ge\lceil rN\rceil\) uses integer cardinalities.
Structural induction, with the inverse witness correspondences and
the ordinary quantifier clauses, proves satisfaction preservation for
the complete finitary formula and every assignment.

If \(\ell=0\), \(N=1,k=1\); there are no witness variables and no
distinct-pair requirements, so the formula reduces to its scope.
All substitutions above are capture-avoiding and retain free parameters.
For rational \(r\), multiplication and integer ceiling compute \(k\)
exactly, so this recursion is an effective construction. With arbitrary
real symbols it is only a semantic construction unless a threshold
comparison/encoding procedure is separately specified.
The construction depends on \(n\); it gives no uniform formula,
size-independent quantifier elimination, complexity bound or complete
axiomatization for unrestricted \(L_{\omega P}\).
\end{proof}
\end{annexitem}

\begin{annexitem}{annex:p15.pq.limit-failures}{thm:p15-pq-limit-failures}{The two source counterexamples and the corrected rate}
\begin{proof}
In the unbiased unary model distribution the number of true instances
is \(K\sim\operatorname{Bin}(n,1/2)\), since all \(n\) ground atoms
are independent fair bits. Complementing all bits pairs counts \(k\)
and \(n-k\) with equal probability. Thus at odd \(n\), the probability
of \(K\ge n/2\) is exactly \(1/2\); at \(n=2m\) it is
\((1+q_m)/2\), where
\[
 q_m=2^{-2m}\binom{2m}{m}
     =\prod_{j=1}^m(1-1/(2j)).
\]
The product identity follows by the factorial ratio for successive
central binomial coefficients. Since \(1-u\le e^{-u}\),
\(q_m\le\exp(-\tfrac12\sum_{j\le m}1/j)\to0\).
The harmonic sum diverges by grouping indices from \(2^a\) to
\(2^{a+1}-1\), each group contributing at least \(1/2\).
This proves the unary limit \(1/2\).

For the binary sentence, both half-threshold assertions on a row hold
exactly when that row has half its entries true. This is impossible at
odd size. At size \(2m\) each row has this probability \(q_m\).
Different rows use disjoint sets of ordered ground relation atoms,
including their diagonal atoms; these sets are independent in the
declared unrestricted model distribution. Consequently
\[
       \mu_{2m}(\eta)=1-(1-q_m)^{2m}.
\]
We prove its limit without assuming the source's quoted Stirling bound.
For \(K\sim\operatorname{Bin}(2m,1/2)\), finite linearity of expectation
and independence give \(\mathbb E K=m\) and
\(\mathbb E(K-m)^2=m/2\): each bit has variance \(1/4\), and
mixed centered terms have expectation zero.
The indicator inequality
\(\mathbf1_{\{|K-m|\ge\sqrt{2m}\}}\le(K-m)^2/(2m)\)
therefore bounds that tail probability by \(1/4\).
At least \(3/4\) of the mass lies on at most
\(2\sqrt{2m}+2\) integer values. The binomial coefficient ratio
\(\binom{2m}{k+1}/\binom{2m}{k}=(2m-k)/(k+1)\)
shows that its maximum mass is the central value \(q_m\).
Hence
\[
 q_m\ge\frac{3}{4(2\sqrt{2m}+2)},\qquad
 (1-q_m)^{2m}\le e^{-2m q_m}\longrightarrow0.
\]
This proves convergence to one along even sizes and zero along odd
sizes, so no convergence law holds for this sentence.

The rate distinction is substantive. Recursion of \(q_m\) gives
\(q_m\le1/\sqrt{m+1}\): the inductive step is the elementary inequality
\((1-1/(2m))^2(m+1)\le m\), valid for \(m\ge1\).
Also \(q_m\le1/2\), and
\(\log(1-u)\ge-2u\) for \(0\le u\le1/2\), as follows by
differentiating \(\log(1-u)+2u\).
Thus
\[
 (1-q_m)^{2m}\ge e^{-4m q_m}
                    \ge e^{-4m/\sqrt{m+1}}.
\]
It cannot have an upper bound \(e^{-cm}\) for any fixed \(c>0\)
at all large \(m\). The source manuscript p.~7 prints a
linear-exponent bound; its convergence conclusion is valid, but that
strong rate is not adopted. The square-root-exponent bound proved
above suffices. Neither counterexample asserts a result for restricted
noncritical fragments, nonuniform tuple measures, undirected graphs,
or dependent ground-atom distributions.
\end{proof}
\end{annexitem}
